%%%%%%%%%%%%%%%%%%%%%%%%%%%%%%%%%%%%%%%%%%%%%%%%%%%%%%%%%%%%
% 2.6 DIOPHANTINE APPROXIMATION
% The Composition of Advanced Mathematics
%%%%%%%%%%%%%%%%%%%%%%%%%%%%%%%%%%%%%%%%%%%%%%%%%%%%%%%%%%%%

\section{Diophantine Approximation}
\label{sec:diophantine-approximation}

\prerequisite{
Elementary Number Theory from \cref{sec:elementary-number-theory}, especially
divisibility, greatest common divisors, rational numbers, and congruences.
Familiarity with sequences, limits, and infinite series from
\cref{sec:real-analysis} is helpful.
}

\learningobjectives{
By the end of this section, the reader should be able to:
\begin{itemize}
    \item explain the central problem of Diophantine approximation;
    \item measure the error in rational approximations to real numbers;
    \item distinguish absolute and normalized approximation error;
    \item use the pigeonhole principle to construct rational approximations;
    \item state and apply Dirichlet's approximation theorem;
    \item explain why irrational numbers admit infinitely many good rational approximations;
    \item construct and analyze continued fractions;
    \item compute convergents of a continued fraction;
    \item use recurrence relations for convergent numerators and denominators;
    \item explain why continued-fraction convergents provide exceptionally good approximations;
    \item recognize the relationship between rational numbers and terminating continued fractions;
    \item recognize periodic continued fractions of quadratic irrationals;
    \item state Lagrange's theorem on periodic continued fractions;
    \item define badly approximable numbers;
    \item explain the significance of the golden ratio in rational approximation;
    \item define the irrationality exponent of a real number;
    \item distinguish algebraic, transcendental, and Liouville-type approximation behavior;
    \item state Liouville's approximation theorem;
    \item explain how unusually strong rational approximation can imply transcendence;
    \item state Roth's theorem conceptually;
    \item interpret approximation using the distance-to-nearest-integer notation;
    \item understand simultaneous Diophantine approximation;
    \item recognize connections with geometry of numbers, dynamical systems,
    transcendental number theory, numerical computation, and number theory.
\end{itemize}
}

%%%%%%%%%%%%%%%%%%%%%%%%%%%%%%%%%%%%%%%%%%%%%%%%%%%%%%%%%%%%
\subsection{The Basic Problem}
\label{subsec:diophantine-approximation-basic-problem}
%%%%%%%%%%%%%%%%%%%%%%%%%%%%%%%%%%%%%%%%%%%%%%%%%%%%%%%%%%%%

Every real number can be approximated by rational numbers.

For example,

\[ \pi\approx\frac{22}{7},\qquad \sqrt2\approx\frac{99}{70}. \]

But not all approximations are equally good.

The central question of Diophantine approximation is:

\begin{center}
\emph{How closely can a real number be approximated by rational numbers?}
\end{center}

The word \emph{Diophantine} refers broadly to questions involving integer or
rational solutions.

\begin{definitionbox}{Diophantine Approximation}
\emph{Diophantine approximation} is the study of approximating real or
complex numbers by rational numbers or, more generally, by algebraic numbers,
subject to arithmetic restrictions.
\end{definitionbox}

A typical approximation has the form

\[ \alpha\approx\frac pq, \]

where

\[ p\in\Z,\qquad q\in\N. \]

%%%%%%%%%%%%%%%%%%%%%%%%%%%%%%%%%%%%%%%%%%%%%%%%%%%%%%%%%%%%
\subsection{Approximation Error}
\label{subsec:diophantine-approximation-error}
%%%%%%%%%%%%%%%%%%%%%%%%%%%%%%%%%%%%%%%%%%%%%%%%%%%%%%%%%%%%

The simplest measure of approximation quality is

\[ \left|\alpha-\frac pq\right|. \]

The smaller this quantity is, the closer \(p/q\) is to \(\alpha\).

For example,

\[ \left|\pi-\frac{22}{7}\right|\approx0.00126449. \]

However, denominator size matters.

A rational number with an enormous denominator can often approximate a real
number very closely simply because many rational numbers are available.

Thus Diophantine approximation asks how small the error can be relative to
the denominator \(q\).

%%%%%%%%%%%%%%%%%%%%%%%%%%%%%%%%%%%%%%%%%%%%%%%%%%%%%%%%%%%%
\subsection{Equivalent Integer Form}
\label{subsec:approximation-integer-form}
%%%%%%%%%%%%%%%%%%%%%%%%%%%%%%%%%%%%%%%%%%%%%%%%%%%%%%%%%%%%

Since

\[ \left|\alpha-\frac pq\right|=\frac{|q\alpha-p|}{q}, \]

the approximation problem can also be expressed by studying

\[ |q\alpha-p|. \]

Thus

\[ \left|\alpha-\frac pq\right|<\frac1{q^2} \]

is equivalent to

\[ |q\alpha-p|<\frac1q. \]

Both forms occur frequently.

%%%%%%%%%%%%%%%%%%%%%%%%%%%%%%%%%%%%%%%%%%%%%%%%%%%%%%%%%%%%
\subsection{Rational Numbers Are Dense}
\label{subsec:rationals-dense-diophantine}
%%%%%%%%%%%%%%%%%%%%%%%%%%%%%%%%%%%%%%%%%%%%%%%%%%%%%%%%%%%%

The rational numbers are dense in \(\R\).

Therefore, for every

\[ \alpha\in\R \]

and every

\[ \varepsilon>0, \]

there exists

\[ \frac pq\in\Q \]

such that

\[ \left|\alpha-\frac pq\right|<\varepsilon. \]

Density guarantees arbitrarily close rational approximations.

Diophantine approximation asks the stronger question:

\[
\boxed{
\text{How does approximation quality depend on the denominator }q?
}
\]

%%%%%%%%%%%%%%%%%%%%%%%%%%%%%%%%%%%%%%%%%%%%%%%%%%%%%%%%%%%%
\subsection{A First Construction}
\label{subsec:first-rational-approximation-construction}
%%%%%%%%%%%%%%%%%%%%%%%%%%%%%%%%%%%%%%%%%%%%%%%%%%%%%%%%%%%%

Fix a positive integer \(q\).

Choose \(p\) to be the integer nearest to \(q\alpha\).

Then

\[ |q\alpha-p|\leq\frac12. \]

Dividing by \(q\),

\[ \left|\alpha-\frac pq\right|\leq\frac1{2q}. \]

This already provides rational approximations.

But much stronger estimates are possible when \(q\) is allowed to vary.

%%%%%%%%%%%%%%%%%%%%%%%%%%%%%%%%%%%%%%%%%%%%%%%%%%%%%%%%%%%%
\subsection{Distance to the Nearest Integer}
\label{subsec:distance-nearest-integer}
%%%%%%%%%%%%%%%%%%%%%%%%%%%%%%%%%%%%%%%%%%%%%%%%%%%%%%%%%%%%

A useful notation is

\[ \|x\|. \]

In this section, when applied to a real number, the notation denotes the
distance from \(x\) to the nearest integer:

\[ \boxed{\|x\|=\min_{m\in\Z}|x-m|.} \]

It is read ``distance of \(x\) to the nearest integer.''

Thus

\[ \|q\alpha\|=\min_{p\in\Z}|q\alpha-p|. \]

\begin{warningbox}
The notation \(\|x\|\) can also denote a norm in other areas of mathematics.
Its meaning must therefore be determined from context.
\end{warningbox}

%%%%%%%%%%%%%%%%%%%%%%%%%%%%%%%%%%%%%%%%%%%%%%%%%%%%%%%%%%%%
\subsection{Visualizing Nearest-Integer Distance}
\label{subsec:nearest-integer-distance-visual}
%%%%%%%%%%%%%%%%%%%%%%%%%%%%%%%%%%%%%%%%%%%%%%%%%%%%%%%%%%%%

\begin{centereddiagram}
\begin{tikzpicture}[x=1.6cm]
\draw[-{Latex[length=2.5mm]}] (-0.3,0)--(5.4,0);

\foreach \x in {0,...,5}{
    \draw (\x,0.09)--(\x,-0.09);
    \node[below=2mm] at (\x,0) {\(\x\)};
}

\coordinate (A) at (3.35,0);
\fill (A) circle (2pt);
\node[above=2mm] at (A) {\(x\)};

\draw[<->] (3,0.42)--(3.35,0.42);
\node[above] at (3.175,0.42) {\(\|x\|\)};
\end{tikzpicture}
\end{centereddiagram}

Here the nearest integer to \(x\) is \(3\), so

\[ \|x\|=|x-3|. \]

%%%%%%%%%%%%%%%%%%%%%%%%%%%%%%%%%%%%%%%%%%%%%%%%%%%%%%%%%%%%
\subsection{The Pigeonhole Principle}
\label{subsec:pigeonhole-diophantine}
%%%%%%%%%%%%%%%%%%%%%%%%%%%%%%%%%%%%%%%%%%%%%%%%%%%%%%%%%%%%

One of the simplest tools for proving rational approximation theorems is the
pigeonhole principle.

Consider the fractional parts

\[ \{0\alpha\},\{1\alpha\},\ldots,\{N\alpha\}. \]

Here

\[ \{x\}=x-\lfloor x\rfloor \]

denotes the fractional part of \(x\).

These \(N+1\) points lie in the interval

\[ [0,1). \]

Divide the interval into \(N\) equal subintervals.

Two of the \(N+1\) fractional parts must lie in the same subinterval.

Their difference is therefore at most

\[ \frac1N. \]

This simple observation produces a powerful theorem.

%%%%%%%%%%%%%%%%%%%%%%%%%%%%%%%%%%%%%%%%%%%%%%%%%%%%%%%%%%%%
\subsection{Dirichlet's Approximation Theorem}
\label{subsec:dirichlet-approximation-theorem}
%%%%%%%%%%%%%%%%%%%%%%%%%%%%%%%%%%%%%%%%%%%%%%%%%%%%%%%%%%%%

\begin{theorem}[Dirichlet's Approximation Theorem]
Let \(\alpha\in\R\) and let \(N\in\N\). Then there exist integers \(p\) and
\(q\) with

\[ 1\leq q\leq N \]

such that

\[ \boxed{\left|\alpha-\frac pq\right|<\frac1{qN}.} \]
\end{theorem}

Equivalently,

\[ \boxed{|q\alpha-p|<\frac1N.} \]

%%%%%%%%%%%%%%%%%%%%%%%%%%%%%%%%%%%%%%%%%%%%%%%%%%%%%%%%%%%%
\subsection{Proof of Dirichlet's Theorem}
\label{subsec:proof-dirichlet-approximation}
%%%%%%%%%%%%%%%%%%%%%%%%%%%%%%%%%%%%%%%%%%%%%%%%%%%%%%%%%%%%

\begin{proof}
Consider the \(N+1\) fractional parts

\[ \{0\alpha\},\{1\alpha\},\ldots,\{N\alpha\}. \]

Partition

\[ [0,1) \]

into \(N\) intervals of length \(1/N\).

By the pigeonhole principle, two fractional parts, say

\[ \{r\alpha\}\qquad\text{and}\qquad\{s\alpha\}, \]

with

\[ 0\leq r<s\leq N, \]

lie in the same interval.

Therefore

\[ |\{s\alpha\}-\{r\alpha\}|<\frac1N. \]

Write

\[ s\alpha=\lfloor s\alpha\rfloor+\{s\alpha\} \]

and similarly for \(r\alpha\).

Then

\[
|(s-r)\alpha-(\lfloor s\alpha\rfloor-\lfloor r\alpha\rfloor)|
<\frac1N.
\]

Set

\[ q=s-r,\qquad p=\lfloor s\alpha\rfloor-\lfloor r\alpha\rfloor. \]

Then

\[ 1\leq q\leq N \]

and

\[ |q\alpha-p|<\frac1N. \]

Dividing by \(q\),

\[ \left|\alpha-\frac pq\right|<\frac1{qN}. \]

This proves the result.
\end{proof}

%%%%%%%%%%%%%%%%%%%%%%%%%%%%%%%%%%%%%%%%%%%%%%%%%%%%%%%%%%%%
\subsection{Infinitely Many Good Approximations}
\label{subsec:infinitely-many-good-rational-approximations}
%%%%%%%%%%%%%%%%%%%%%%%%%%%%%%%%%%%%%%%%%%%%%%%%%%%%%%%%%%%%

If \(\alpha\) is irrational, Dirichlet's theorem implies that there are
infinitely many rational numbers \(p/q\) satisfying

\[ \boxed{\left|\alpha-\frac pq\right|<\frac1{q^2}.} \]

Indeed, apply Dirichlet's theorem for arbitrarily large \(N\).

Since

\[ q\leq N, \]

we have

\[ \frac1{qN}\leq\frac1{q^2}. \]

Thus irrational numbers always admit infinitely many approximations with
quadratic denominator error.

%%%%%%%%%%%%%%%%%%%%%%%%%%%%%%%%%%%%%%%%%%%%%%%%%%%%%%%%%%%%
\subsection{Why the Exponent \(2\) Matters}
\label{subsec:approximation-exponent-two}
%%%%%%%%%%%%%%%%%%%%%%%%%%%%%%%%%%%%%%%%%%%%%%%%%%%%%%%%%%%%

The estimate

\[ \left|\alpha-\frac pq\right|<\frac1{q^2} \]

is much stronger than the elementary estimate

\[ O\left(\frac1q\right). \]

The exponent \(2\) becomes a fundamental dividing line.

Many questions in Diophantine approximation ask whether an inequality of the
form

\[ \left|\alpha-\frac pq\right|<\frac1{q^\mu} \]

has infinitely many rational solutions.

The parameter

\[ \mu \]

is lowercase Greek \emph{mu}, pronounced ``mew.''

%%%%%%%%%%%%%%%%%%%%%%%%%%%%%%%%%%%%%%%%%%%%%%%%%%%%%%%%%%%%
\subsection{Continued Fractions}
\label{subsec:continued-fractions}
%%%%%%%%%%%%%%%%%%%%%%%%%%%%%%%%%%%%%%%%%%%%%%%%%%%%%%%%%%%%

Continued fractions provide a systematic method for constructing exceptionally
good rational approximations.

A finite continued fraction has the form

\[
a_0+\cfrac{1}{
a_1+\cfrac{1}{
a_2+\cfrac{1}{
\ddots+\cfrac1{a_n}}}}.
\]

This is abbreviated

\[ [a_0;a_1,a_2,\ldots,a_n]. \]

The notation is read

\begin{center}
``\(a_0\); \(a_1,a_2,\ldots,a_n\) as a continued fraction.''
\end{center}

Usually

\[ a_0\in\Z,\qquad a_1,a_2,\ldots\in\N. \]

%%%%%%%%%%%%%%%%%%%%%%%%%%%%%%%%%%%%%%%%%%%%%%%%%%%%%%%%%%%%
\subsection{Simple Continued Fractions}
\label{subsec:simple-continued-fractions}
%%%%%%%%%%%%%%%%%%%%%%%%%%%%%%%%%%%%%%%%%%%%%%%%%%%%%%%%%%%%

\begin{definitionbox}{Simple Continued Fraction}
A \emph{simple continued fraction} is an expression

\[ [a_0;a_1,a_2,\ldots] \]

where

\[ a_0\in\Z \]

and

\[ a_n\in\N\qquad(n\geq1). \]
\end{definitionbox}

Finite simple continued fractions represent rational numbers.

Infinite simple continued fractions represent irrational numbers.

%%%%%%%%%%%%%%%%%%%%%%%%%%%%%%%%%%%%%%%%%%%%%%%%%%%%%%%%%%%%
\subsection{Constructing a Continued Fraction}
\label{subsec:constructing-continued-fraction}
%%%%%%%%%%%%%%%%%%%%%%%%%%%%%%%%%%%%%%%%%%%%%%%%%%%%%%%%%%%%

Let

\[ \alpha\in\R. \]

Set

\[ a_0=\lfloor\alpha\rfloor. \]

Then

\[ \alpha=a_0+r_0,\qquad0\leq r_0<1. \]

If \(r_0\neq0\), invert the fractional part:

\[ \frac1{r_0}=a_1+r_1. \]

Continue recursively.

This produces

\[ \alpha=[a_0;a_1,a_2,\ldots]. \]

%%%%%%%%%%%%%%%%%%%%%%%%%%%%%%%%%%%%%%%%%%%%%%%%%%%%%%%%%%%%
\subsection{Continued Fraction of a Rational Number}
\label{subsec:rational-continued-fraction}
%%%%%%%%%%%%%%%%%%%%%%%%%%%%%%%%%%%%%%%%%%%%%%%%%%%%%%%%%%%%

Consider

\[ \frac{43}{13}. \]

Using the Euclidean Algorithm,

\[
43=3(13)+4,
\]

\[
13=3(4)+1,
\]

\[
4=4(1).
\]

Therefore

\[ \boxed{\frac{43}{13}=[3;3,4].} \]

The Euclidean Algorithm and finite continued fractions are therefore two
forms of the same arithmetic process.

%%%%%%%%%%%%%%%%%%%%%%%%%%%%%%%%%%%%%%%%%%%%%%%%%%%%%%%%%%%%
\subsection{Worked Example: Continued Fraction Expansion}
\label{subsec:worked-continued-fraction-expansion}
%%%%%%%%%%%%%%%%%%%%%%%%%%%%%%%%%%%%%%%%%%%%%%%%%%%%%%%%%%%%

\begin{workedexamplebox}{Expanding \(\frac{17}{7}\)}

Write

\[ \frac{17}{7} \]

as a simple continued fraction.

\begin{tabularx}{\textwidth}{@{}>{\raggedright\arraybackslash}p{0.43\textwidth}X@{}}
Divide \(17\) by \(7\). & \(17=2(7)+3\)\\
Invert the remainder ratio. & \(\frac73=2+\frac13\)\\
Invert again. & \(\frac31=3\)\\
Collect the quotients. & \([2;2,3]\)
\end{tabularx}

\begin{answerbox}
\[ \boxed{\frac{17}{7}=[2;2,3].} \]
\end{answerbox}
\end{workedexamplebox}

%%%%%%%%%%%%%%%%%%%%%%%%%%%%%%%%%%%%%%%%%%%%%%%%%%%%%%%%%%%%
\subsection{Infinite Continued Fractions}
\label{subsec:infinite-continued-fractions}
%%%%%%%%%%%%%%%%%%%%%%%%%%%%%%%%%%%%%%%%%%%%%%%%%%%%%%%%%%%%

An irrational number has an infinite continued-fraction expansion:

\[ \alpha=[a_0;a_1,a_2,a_3,\ldots]. \]

Truncating the expansion produces rational approximations:

\[
[a_0],\qquad
[a_0;a_1],\qquad
[a_0;a_1,a_2],\qquad\ldots
\]

These are called convergents.

%%%%%%%%%%%%%%%%%%%%%%%%%%%%%%%%%%%%%%%%%%%%%%%%%%%%%%%%%%%%
\subsection{Convergents}
\label{subsec:continued-fraction-convergents}
%%%%%%%%%%%%%%%%%%%%%%%%%%%%%%%%%%%%%%%%%%%%%%%%%%%%%%%%%%%%

\begin{definitionbox}{Convergent}
The \(n\)-th \emph{convergent} of

\[ \alpha=[a_0;a_1,a_2,\ldots] \]

is the rational number

\[ \frac{p_n}{q_n}=[a_0;a_1,\ldots,a_n]. \]
\end{definitionbox}

The convergents approach \(\alpha\) extremely efficiently.

%%%%%%%%%%%%%%%%%%%%%%%%%%%%%%%%%%%%%%%%%%%%%%%%%%%%%%%%%%%%
\subsection{Recurrence Relations for Convergents}
\label{subsec:convergent-recurrences}
%%%%%%%%%%%%%%%%%%%%%%%%%%%%%%%%%%%%%%%%%%%%%%%%%%%%%%%%%%%%

Define

\[
p_{-2}=0,\qquad p_{-1}=1,
\]

and

\[
q_{-2}=1,\qquad q_{-1}=0.
\]

Then

\[ \boxed{p_n=a_np_{n-1}+p_{n-2}} \]

and

\[ \boxed{q_n=a_nq_{n-1}+q_{n-2}}. \]

Thus the convergents can be generated recursively without repeatedly
evaluating nested fractions.

%%%%%%%%%%%%%%%%%%%%%%%%%%%%%%%%%%%%%%%%%%%%%%%%%%%%%%%%%%%%
\subsection{Determinant Identity}
\label{subsec:continued-fraction-determinant}
%%%%%%%%%%%%%%%%%%%%%%%%%%%%%%%%%%%%%%%%%%%%%%%%%%%%%%%%%%%%

Consecutive convergents satisfy

\[ \boxed{p_nq_{n-1}-p_{n-1}q_n=(-1)^{n-1}.} \]

Consequently,

\[ \gcd(p_n,q_n)=1. \]

Thus every convergent is automatically written in lowest terms.

The identity also gives

\[
\left|
\frac{p_n}{q_n}-\frac{p_{n-1}}{q_{n-1}}
\right|
=
\frac1{q_nq_{n-1}}.
\]

%%%%%%%%%%%%%%%%%%%%%%%%%%%%%%%%%%%%%%%%%%%%%%%%%%%%%%%%%%%%
\subsection{Convergents Alternate Around the Limit}
\label{subsec:convergents-alternate}
%%%%%%%%%%%%%%%%%%%%%%%%%%%%%%%%%%%%%%%%%%%%%%%%%%%%%%%%%%%%

For an irrational number \(\alpha\), its convergents alternate around
\(\alpha\).

Schematically,

\[
\frac{p_0}{q_0}
<
\frac{p_2}{q_2}
<
\cdots
<
\alpha
<
\cdots
<
\frac{p_3}{q_3}
<
\frac{p_1}{q_1},
\]

with orientation depending on the initial expansion.

\begin{centereddiagram}
\begin{tikzpicture}[x=1.25cm]
\draw[-{Latex[length=2.5mm]}] (0,0)--(8.2,0);

\draw (4,0.13)--(4,-0.13);
\node[above=3mm] at (4,0) {\(\alpha\)};

\fill (1.1,0) circle (2pt);
\fill (2.6,0) circle (2pt);
\fill (3.45,0) circle (2pt);
\fill (4.55,0) circle (2pt);
\fill (5.5,0) circle (2pt);
\fill (7,0) circle (2pt);

\node[below=2mm] at (1.1,0) {\(\frac{p_0}{q_0}\)};
\node[below=2mm] at (2.6,0) {\(\frac{p_2}{q_2}\)};
\node[below=2mm] at (3.45,0) {\(\frac{p_4}{q_4}\)};
\node[below=2mm] at (4.55,0) {\(\frac{p_5}{q_5}\)};
\node[below=2mm] at (5.5,0) {\(\frac{p_3}{q_3}\)};
\node[below=2mm] at (7,0) {\(\frac{p_1}{q_1}\)};
\end{tikzpicture}
\end{centereddiagram}

The approximations squeeze toward \(\alpha\) from both sides.

%%%%%%%%%%%%%%%%%%%%%%%%%%%%%%%%%%%%%%%%%%%%%%%%%%%%%%%%%%%%
\subsection{Quality of Continued-Fraction Approximations}
\label{subsec:continued-fraction-quality}
%%%%%%%%%%%%%%%%%%%%%%%%%%%%%%%%%%%%%%%%%%%%%%%%%%%%%%%%%%%%

A fundamental estimate is

\begin{theorem}
If \(p_n/q_n\) is a convergent of an irrational number \(\alpha\), then

\[ \boxed{\left|\alpha-\frac{p_n}{q_n}\right|<\frac1{q_nq_{n+1}}<\frac1{q_n^2}.} \]
\end{theorem}

Thus every convergent gives an approximation at least as strong as the
quadratic scale predicted by Dirichlet's theorem.

%%%%%%%%%%%%%%%%%%%%%%%%%%%%%%%%%%%%%%%%%%%%%%%%%%%%%%%%%%%%
\subsection{Best Approximation Property}
\label{subsec:continued-fraction-best-approximation}
%%%%%%%%%%%%%%%%%%%%%%%%%%%%%%%%%%%%%%%%%%%%%%%%%%%%%%%%%%%%

Continued-fraction convergents are not merely convenient approximations.

They are essentially the best possible rational approximations for their
denominator sizes.

One standard form of the result states that if

\[ 0<q<q_{n+1}, \]

then no rational \(p/q\) improves substantially upon the approximation
provided by the \(n\)-th convergent.

\begin{conceptbox}
Continued fractions organize rational approximations in order of arithmetic
efficiency: very small error using relatively small denominators.
\end{conceptbox}

%%%%%%%%%%%%%%%%%%%%%%%%%%%%%%%%%%%%%%%%%%%%%%%%%%%%%%%%%%%%
\subsection{Continued Fraction of \(\sqrt2\)}
\label{subsec:continued-fraction-sqrt2}
%%%%%%%%%%%%%%%%%%%%%%%%%%%%%%%%%%%%%%%%%%%%%%%%%%%%%%%%%%%%

Let

\[ \alpha=\sqrt2. \]

Since

\[ 1<\sqrt2<2, \]

we begin with

\[ a_0=1. \]

Now

\[
\frac1{\sqrt2-1}
=
\sqrt2+1
=
2+(\sqrt2-1).
\]

The same fractional part reappears.

Therefore

\[ \boxed{\sqrt2=[1;2,2,2,2,\ldots].} \]

This is often written

\[ \boxed{\sqrt2=[1;\overline{2}].} \]

The bar indicates that the block repeats indefinitely.

%%%%%%%%%%%%%%%%%%%%%%%%%%%%%%%%%%%%%%%%%%%%%%%%%%%%%%%%%%%%
\subsection{Convergents of \(\sqrt2\)}
\label{subsec:sqrt2-convergents}
%%%%%%%%%%%%%%%%%%%%%%%%%%%%%%%%%%%%%%%%%%%%%%%%%%%%%%%%%%%%

The first convergents are

\[
1,\qquad
\frac32,\qquad
\frac75,\qquad
\frac{17}{12},\qquad
\frac{41}{29},\qquad
\frac{99}{70}.
\]

Compare:

\[
\sqrt2\approx1.414213562\ldots
\]

while

\[ \frac{99}{70}\approx1.414285714. \]

The approximation is already excellent despite the modest denominator.

%%%%%%%%%%%%%%%%%%%%%%%%%%%%%%%%%%%%%%%%%%%%%%%%%%%%%%%%%%%%
\subsection{Worked Example: Approximating \(\sqrt2\)}
\label{subsec:worked-sqrt2-convergent}
%%%%%%%%%%%%%%%%%%%%%%%%%%%%%%%%%%%%%%%%%%%%%%%%%%%%%%%%%%%%

\begin{workedexamplebox}{A Convergent of \(\sqrt2\)}

Use the continued fraction

\[ \sqrt2=[1;\overline2] \]

to compute the fourth convergent.

\begin{tabularx}{\textwidth}{@{}>{\raggedright\arraybackslash}p{0.43\textwidth}X@{}}
Start with the truncation. & \([1;2,2,2]\)\\
Evaluate from the bottom. & \(2+\frac12=\frac52\)\\
Continue upward. & \(2+\frac{2}{5}=\frac{12}{5}\)\\
Invert and add \(1\). & \(1+\frac5{12}=\frac{17}{12}\)
\end{tabularx}

\begin{answerbox}
\[ \boxed{\sqrt2\approx\frac{17}{12}.} \]
\end{answerbox}
\end{workedexamplebox}

%%%%%%%%%%%%%%%%%%%%%%%%%%%%%%%%%%%%%%%%%%%%%%%%%%%%%%%%%%%%
\subsection{The Golden Ratio}
\label{subsec:golden-ratio-diophantine}
%%%%%%%%%%%%%%%%%%%%%%%%%%%%%%%%%%%%%%%%%%%%%%%%%%%%%%%%%%%%

The golden ratio is

\[ \varphi=\frac{1+\sqrt5}{2}. \]

Here \(\varphi\) is lowercase Greek \emph{phi}.

It satisfies

\[ \varphi=1+\frac1\varphi. \]

Therefore its continued fraction repeats forever:

\[ \boxed{\varphi=[1;1,1,1,1,\ldots]=[1;\overline1].} \]

Its convergents are ratios of consecutive Fibonacci numbers:

\[
1,\quad
2,\quad
\frac32,\quad
\frac53,\quad
\frac85,\quad
\frac{13}{8},\quad
\frac{21}{13},\ldots
\]

%%%%%%%%%%%%%%%%%%%%%%%%%%%%%%%%%%%%%%%%%%%%%%%%%%%%%%%%%%%%
\subsection{Fibonacci Convergents}
\label{subsec:fibonacci-golden-convergents}
%%%%%%%%%%%%%%%%%%%%%%%%%%%%%%%%%%%%%%%%%%%%%%%%%%%%%%%%%%%%

Let the Fibonacci numbers satisfy

\[ F_{n+1}=F_n+F_{n-1}. \]

Because every partial quotient of the golden ratio equals \(1\), the
continued-fraction recurrence becomes exactly the Fibonacci recurrence.

Thus

\[ \boxed{\frac{F_{n+1}}{F_n}\to\varphi.} \]

This gives a beautiful connection among:

\[
\text{continued fractions}
\longleftrightarrow
\text{Fibonacci numbers}
\longleftrightarrow
\text{the golden ratio}.
\]

%%%%%%%%%%%%%%%%%%%%%%%%%%%%%%%%%%%%%%%%%%%%%%%%%%%%%%%%%%%%
\subsection{Badly Approximable Numbers}
\label{subsec:badly-approximable-numbers}
%%%%%%%%%%%%%%%%%%%%%%%%%%%%%%%%%%%%%%%%%%%%%%%%%%%%%%%%%%%%

Some irrational numbers resist unusually close rational approximation.

\begin{definitionbox}{Badly Approximable Number}
An irrational number \(\alpha\) is \emph{badly approximable} if there exists
a constant \(c>0\) such that

\[ \boxed{\left|\alpha-\frac pq\right|>\frac{c}{q^2}} \]

for every rational number \(p/q\).
\end{definitionbox}

The terminology may sound misleading.

A badly approximable number still has infinitely many good rational
approximations.

It simply cannot be approximated substantially better than the natural
\(1/q^2\) scale.

%%%%%%%%%%%%%%%%%%%%%%%%%%%%%%%%%%%%%%%%%%%%%%%%%%%%%%%%%%%%
\subsection{Continued-Fraction Characterization}
\label{subsec:badly-approximable-cf-characterization}
%%%%%%%%%%%%%%%%%%%%%%%%%%%%%%%%%%%%%%%%%%%%%%%%%%%%%%%%%%%%

A fundamental theorem states:

\begin{theorem}
An irrational real number is badly approximable if and only if the partial
quotients in its continued fraction are bounded.
\end{theorem}

Thus the size of the coefficients

\[ a_1,a_2,a_3,\ldots \]

directly controls rational approximation quality.

Large partial quotients correspond to exceptionally good convergents.

%%%%%%%%%%%%%%%%%%%%%%%%%%%%%%%%%%%%%%%%%%%%%%%%%%%%%%%%%%%%
\subsection{The Golden Ratio as an Extremal Number}
\label{subsec:golden-ratio-hardest-approximate}
%%%%%%%%%%%%%%%%%%%%%%%%%%%%%%%%%%%%%%%%%%%%%%%%%%%%%%%%%%%%

The golden ratio has the continued fraction

\[ [1;\overline1]. \]

Its partial quotients are as small as possible.

Consequently, the denominators of its convergents grow relatively slowly,
and its rational approximations improve relatively slowly.

In a precise sense related to Hurwitz's theorem, the golden ratio is the
classical extremal example of a number that is difficult to approximate too
closely by rationals.

%%%%%%%%%%%%%%%%%%%%%%%%%%%%%%%%%%%%%%%%%%%%%%%%%%%%%%%%%%%%
\subsection{Hurwitz's Approximation Theorem}
\label{subsec:hurwitz-theorem}
%%%%%%%%%%%%%%%%%%%%%%%%%%%%%%%%%%%%%%%%%%%%%%%%%%%%%%%%%%%%

Dirichlet's theorem gives the constant \(1\) in

\[ \left|\alpha-\frac pq\right|<\frac1{q^2}. \]

Hurwitz improved the universal constant.

\begin{theorem}[Hurwitz]
For every irrational real number \(\alpha\), there exist infinitely many
rational numbers \(p/q\) such that

\[ \boxed{\left|\alpha-\frac pq\right|<\frac1{\sqrt5\,q^2}.} \]
\end{theorem}

The constant

\[ \frac1{\sqrt5} \]

cannot be replaced by a smaller universal constant.

The golden ratio explains why this constant is optimal.

%%%%%%%%%%%%%%%%%%%%%%%%%%%%%%%%%%%%%%%%%%%%%%%%%%%%%%%%%%%%
\subsection{Quadratic Irrational Numbers}
\label{subsec:quadratic-irrationals-diophantine}
%%%%%%%%%%%%%%%%%%%%%%%%%%%%%%%%%%%%%%%%%%%%%%%%%%%%%%%%%%%%

\begin{definitionbox}{Quadratic Irrational}
A \emph{quadratic irrational} is an irrational real number satisfying a
quadratic equation

\[ ax^2+bx+c=0 \]

with

\[ a,b,c\in\Z,\qquad a\neq0. \]
\end{definitionbox}

Examples include

\[ \sqrt2,\qquad\sqrt3,\qquad\frac{1+\sqrt5}{2}. \]

Quadratic irrationals have a remarkable continued-fraction structure.

%%%%%%%%%%%%%%%%%%%%%%%%%%%%%%%%%%%%%%%%%%%%%%%%%%%%%%%%%%%%
\subsection{Lagrange's Theorem}
\label{subsec:lagrange-periodic-continued-fractions}
%%%%%%%%%%%%%%%%%%%%%%%%%%%%%%%%%%%%%%%%%%%%%%%%%%%%%%%%%%%%

\begin{theorem}[Lagrange]
A real number has an eventually periodic simple continued fraction if and
only if it is a quadratic irrational.
\end{theorem}

Thus

\[
\boxed{
\text{quadratic irrational}
\iff
\text{eventually periodic continued fraction}.
}
\]

For example,

\[ \sqrt2=[1;\overline2]. \]

This theorem gives an exact arithmetic characterization of quadratic
irrationality.

%%%%%%%%%%%%%%%%%%%%%%%%%%%%%%%%%%%%%%%%%%%%%%%%%%%%%%%%%%%%
\subsection{Pell Equations}
\label{subsec:pell-equations-diophantine-approximation}
%%%%%%%%%%%%%%%%%%%%%%%%%%%%%%%%%%%%%%%%%%%%%%%%%%%%%%%%%%%%

Continued fractions are closely related to equations of the form

\[ x^2-Dy^2=1, \]

where \(D\) is a positive nonsquare integer.

These are called \emph{Pell equations}.

If

\[ \frac xy\approx\sqrt D, \]

then

\[ x^2-Dy^2 \]

may be small.

Continued-fraction convergents of \(\sqrt D\) systematically produce
solutions to Pell equations.

This creates a bridge between Diophantine approximation and Diophantine
equations.

%%%%%%%%%%%%%%%%%%%%%%%%%%%%%%%%%%%%%%%%%%%%%%%%%%%%%%%%%%%%
\subsection{Example: Pell's Equation for \(D=2\)}
\label{subsec:pell-d2}
%%%%%%%%%%%%%%%%%%%%%%%%%%%%%%%%%%%%%%%%%%%%%%%%%%%%%%%%%%%%

The convergent

\[ \frac32 \]

of \(\sqrt2\) satisfies

\[ 3^2-2(2^2)=1. \]

Likewise,

\[ \frac{17}{12} \]

satisfies

\[ 17^2-2(12^2)=1. \]

Thus

\[ (x,y)=(3,2),\qquad(17,12) \]

are solutions of

\[ x^2-2y^2=1. \]

%%%%%%%%%%%%%%%%%%%%%%%%%%%%%%%%%%%%%%%%%%%%%%%%%%%%%%%%%%%%
\subsection{Irrationality Measures}
\label{subsec:irrationality-measures}
%%%%%%%%%%%%%%%%%%%%%%%%%%%%%%%%%%%%%%%%%%%%%%%%%%%%%%%%%%%%

Suppose \(\alpha\) is irrational.

We can ask how large an exponent \(\mu\) can occur infinitely often in

\[ \left|\alpha-\frac pq\right|<\frac1{q^\mu}. \]

This leads to the irrationality exponent.

%%%%%%%%%%%%%%%%%%%%%%%%%%%%%%%%%%%%%%%%%%%%%%%%%%%%%%%%%%%%
\subsection{Irrationality Exponent}
\label{subsec:irrationality-exponent}
%%%%%%%%%%%%%%%%%%%%%%%%%%%%%%%%%%%%%%%%%%%%%%%%%%%%%%%%%%%%

\begin{definitionbox}{Irrationality Exponent}
For an irrational real number \(\alpha\), the \emph{irrationality exponent}
\(\mu(\alpha)\) is the supremum of the real numbers \(\mu\) for which

\[ 0<\left|\alpha-\frac pq\right|<\frac1{q^\mu} \]

has infinitely many rational solutions \(p/q\).
\end{definitionbox}

It is read

\begin{center}
``mu of alpha.''
\end{center}

Dirichlet's theorem implies

\[ \boxed{\mu(\alpha)\geq2} \]

for every irrational \(\alpha\).

%%%%%%%%%%%%%%%%%%%%%%%%%%%%%%%%%%%%%%%%%%%%%%%%%%%%%%%%%%%%
\subsection{Meaning of the Irrationality Exponent}
\label{subsec:meaning-irrationality-exponent}
%%%%%%%%%%%%%%%%%%%%%%%%%%%%%%%%%%%%%%%%%%%%%%%%%%%%%%%%%%%%

A larger value of

\[ \mu(\alpha) \]

means that \(\alpha\) admits exceptionally accurate rational approximations.

Very roughly:

\[
\begin{array}{c|c}
\mu(\alpha) & \text{Approximation behavior}\\
\hline
2 & \text{ordinary irrational approximation}\\
>2 & \text{unusually strong approximation}\\
\infty & \text{extremely strong approximation}
\end{array}
\]

The irrationality exponent therefore quantifies how ``rational-like'' an
irrational number can appear at selected denominator scales.

%%%%%%%%%%%%%%%%%%%%%%%%%%%%%%%%%%%%%%%%%%%%%%%%%%%%%%%%%%%%
\subsection{Algebraic Numbers and Approximation}
\label{subsec:algebraic-number-approximation}
%%%%%%%%%%%%%%%%%%%%%%%%%%%%%%%%%%%%%%%%%%%%%%%%%%%%%%%%%%%%

Algebraic irrational numbers cannot be approximated arbitrarily well by
rationals.

The first major theorem establishing this phenomenon was proved by
Liouville.

%%%%%%%%%%%%%%%%%%%%%%%%%%%%%%%%%%%%%%%%%%%%%%%%%%%%%%%%%%%%
\subsection{Liouville's Approximation Theorem}
\label{subsec:liouville-approximation-theorem}
%%%%%%%%%%%%%%%%%%%%%%%%%%%%%%%%%%%%%%%%%%%%%%%%%%%%%%%%%%%%

\begin{theorem}[Liouville]
Let \(\alpha\) be an algebraic irrational number of degree \(d\geq2\).
Then there exists a constant \(C(\alpha)>0\) such that

\[ \boxed{\left|\alpha-\frac pq\right|>\frac{C(\alpha)}{q^d}} \]

for every rational \(p/q\) with \(q>0\).
\end{theorem}

The notation

\[ C(\alpha) \]

means a positive constant depending on \(\alpha\).

Liouville's theorem places a fundamental restriction on how well algebraic
numbers can be approximated by rationals.

%%%%%%%%%%%%%%%%%%%%%%%%%%%%%%%%%%%%%%%%%%%%%%%%%%%%%%%%%%%%
\subsection{Why Liouville's Theorem Suggests Transcendence}
\label{subsec:liouville-transcendence-idea}
%%%%%%%%%%%%%%%%%%%%%%%%%%%%%%%%%%%%%%%%%%%%%%%%%%%%%%%%%%%%

Suppose a real number can be approximated so well that for arbitrarily large
exponents \(m\),

\[ 0<\left|\alpha-\frac pq\right|<\frac1{q^m} \]

has rational solutions.

No algebraic number of fixed degree can behave this way.

Therefore such a number must be transcendental.

This was one of the earliest systematic methods for proving transcendence.

%%%%%%%%%%%%%%%%%%%%%%%%%%%%%%%%%%%%%%%%%%%%%%%%%%%%%%%%%%%%
\subsection{Liouville Numbers}
\label{subsec:liouville-numbers}
%%%%%%%%%%%%%%%%%%%%%%%%%%%%%%%%%%%%%%%%%%%%%%%%%%%%%%%%%%%%

\begin{definitionbox}{Liouville Number}
A real number \(\alpha\) is a \emph{Liouville number} if, for every positive
integer \(m\), there exist integers \(p\) and \(q>1\) such that

\[ 0<\left|\alpha-\frac pq\right|<\frac1{q^m}. \]
\end{definitionbox}

Every Liouville number is transcendental.

Equivalently, a Liouville number has

\[ \boxed{\mu(\alpha)=\infty.} \]

%%%%%%%%%%%%%%%%%%%%%%%%%%%%%%%%%%%%%%%%%%%%%%%%%%%%%%%%%%%%
\subsection{Liouville's Constant}
\label{subsec:liouville-constant}
%%%%%%%%%%%%%%%%%%%%%%%%%%%%%%%%%%%%%%%%%%%%%%%%%%%%%%%%%%%%

A famous example is

\[ L=\sum_{n=1}^{\infty}10^{-n!}. \]

Thus

\[
L
=
0.110001000000000000000001\ldots
\]

The nonzero decimal digits occur at factorial positions.

Truncating the series after the \(m\)-th term produces a rational number with
denominator

\[ 10^{m!}. \]

The remaining tail is extraordinarily small.

This gives rational approximations far better than any fixed power of the
denominator.

Hence \(L\) is transcendental.

%%%%%%%%%%%%%%%%%%%%%%%%%%%%%%%%%%%%%%%%%%%%%%%%%%%%%%%%%%%%
\subsection{Worked Example: Liouville Approximation}
\label{subsec:worked-liouville-approximation}
%%%%%%%%%%%%%%%%%%%%%%%%%%%%%%%%%%%%%%%%%%%%%%%%%%%%%%%%%%%%

\begin{workedexamplebox}{A Truncation of Liouville's Constant}

Use the first three terms of

\[ L=\sum_{n=1}^{\infty}10^{-n!}. \]

\begin{tabularx}{\textwidth}{@{}>{\raggedright\arraybackslash}p{0.43\textwidth}X@{}}
Write the first three terms. &
\(10^{-1!}+10^{-2!}+10^{-3!}\)\\
Evaluate the exponents. &
\(10^{-1}+10^{-2}+10^{-6}\)\\
Add. &
\(0.1+0.01+0.000001\)\\
Obtain the rational approximation. &
\(0.110001=\frac{110001}{10^6}\)
\end{tabularx}

\begin{answerbox}
\[ \boxed{L\approx\frac{110001}{10^6}.} \]
\end{answerbox}
\end{workedexamplebox}

%%%%%%%%%%%%%%%%%%%%%%%%%%%%%%%%%%%%%%%%%%%%%%%%%%%%%%%%%%%%
\subsection{Roth's Theorem}
\label{subsec:roth-theorem}
%%%%%%%%%%%%%%%%%%%%%%%%%%%%%%%%%%%%%%%%%%%%%%%%%%%%%%%%%%%%

Liouville's theorem was later strengthened dramatically.

\begin{theorem}[Roth]
Let \(\alpha\) be an algebraic irrational number. For every
\(\varepsilon>0\), the inequality

\[ \boxed{\left|\alpha-\frac pq\right|<\frac1{q^{2+\varepsilon}}} \]

has only finitely many rational solutions \(p/q\).
\end{theorem}

Therefore every algebraic irrational number satisfies

\[ \boxed{\mu(\alpha)=2.} \]

This is optimal because Dirichlet's theorem guarantees exponent \(2\)
infinitely often.

%%%%%%%%%%%%%%%%%%%%%%%%%%%%%%%%%%%%%%%%%%%%%%%%%%%%%%%%%%%%
\subsection{The Approximation Threshold}
\label{subsec:approximation-threshold}
%%%%%%%%%%%%%%%%%%%%%%%%%%%%%%%%%%%%%%%%%%%%%%%%%%%%%%%%%%%%

Dirichlet and Roth together create a striking picture for algebraic
irrational numbers:

\[
\boxed{
\begin{array}{c}
\text{Exponent }2\text{ occurs infinitely often},\\[1mm]
\text{every exponent }2+\varepsilon\text{ occurs only finitely often}.
\end{array}
}
\]

Thus \(2\) is the exact rational-approximation exponent for algebraic
irrational numbers.

%%%%%%%%%%%%%%%%%%%%%%%%%%%%%%%%%%%%%%%%%%%%%%%%%%%%%%%%%%%%
\subsection{Transcendental Numbers Need Not Be Liouville}
\label{subsec:transcendental-not-liouville}
%%%%%%%%%%%%%%%%%%%%%%%%%%%%%%%%%%%%%%%%%%%%%%%%%%%%%%%%%%%%

It is important not to conclude that every transcendental number is
exceptionally well approximable.

Many transcendental numbers have irrationality exponent exactly \(2\).

Thus

\[
\text{Liouville}
\Longrightarrow
\text{transcendental},
\]

but the converse is false.

The classification of approximation behavior among transcendental numbers
is much richer.

%%%%%%%%%%%%%%%%%%%%%%%%%%%%%%%%%%%%%%%%%%%%%%%%%%%%%%%%%%%%
\subsection{Approximation of \(\pi\)}
\label{subsec:pi-rational-approximation}
%%%%%%%%%%%%%%%%%%%%%%%%%%%%%%%%%%%%%%%%%%%%%%%%%%%%%%%%%%%%

The familiar approximation

\[ \frac{22}{7} \]

is good, but continued fractions produce substantially better approximations.

The continued fraction of \(\pi\) begins

\[ \pi=[3;7,15,1,292,\ldots]. \]

Its convergents include

\[
3,\qquad
\frac{22}{7},\qquad
\frac{333}{106},\qquad
\frac{355}{113}.
\]

The approximation

\[ \boxed{\frac{355}{113}\approx3.14159292035} \]

is remarkably accurate for such a small denominator.

%%%%%%%%%%%%%%%%%%%%%%%%%%%%%%%%%%%%%%%%%%%%%%%%%%%%%%%%%%%%
\subsection{Why \(355/113\) Is So Accurate}
\label{subsec:pi-355-113}
%%%%%%%%%%%%%%%%%%%%%%%%%%%%%%%%%%%%%%%%%%%%%%%%%%%%%%%%%%%%

The next partial quotient after

\[ \frac{355}{113} \]

in the continued fraction of \(\pi\) is unusually large.

Large partial quotients produce exceptionally accurate preceding
convergents.

Thus the continued fraction itself explains why

\[ \frac{355}{113} \]

is unexpectedly good.

%%%%%%%%%%%%%%%%%%%%%%%%%%%%%%%%%%%%%%%%%%%%%%%%%%%%%%%%%%%%
\subsection{Approximation on the Number Line}
\label{subsec:rational-approximation-number-line}
%%%%%%%%%%%%%%%%%%%%%%%%%%%%%%%%%%%%%%%%%%%%%%%%%%%%%%%%%%%%

\begin{centereddiagram}
\begin{tikzpicture}[x=3.2cm]
\draw[-{Latex[length=2.5mm]}] (0,0)--(3.5,0);

\coordinate (P) at (1.95,0);
\coordinate (A) at (2.12,0);

\draw (P) ++(0,0.1)--++(0,-0.2);
\draw (A) ++(0,0.1)--++(0,-0.2);

\node[below=2mm] at (P) {\(\frac pq\)};
\node[below=2mm] at (A) {\(\alpha\)};

\draw[<->] ($(P)+(0,0.42)$)--($(A)+(0,0.42)$);
\node[above] at ($(P)!0.5!(A)+(0,0.42)$)
{\(\left|\alpha-\frac pq\right|\)};
\end{tikzpicture}
\end{centereddiagram}

Diophantine approximation seeks rational points that approach \(\alpha\)
rapidly while keeping \(q\) relatively small.

%%%%%%%%%%%%%%%%%%%%%%%%%%%%%%%%%%%%%%%%%%%%%%%%%%%%%%%%%%%%
\subsection{Simultaneous Diophantine Approximation}
\label{subsec:simultaneous-diophantine-approximation}
%%%%%%%%%%%%%%%%%%%%%%%%%%%%%%%%%%%%%%%%%%%%%%%%%%%%%%%%%%%%

The approximation problem can involve several real numbers at once.

Given

\[ \alpha_1,\ldots,\alpha_n, \]

we may seek a common denominator \(q\) and integers

\[ p_1,\ldots,p_n \]

such that

\[ q\alpha_i\approx p_i \]

for every \(i\).

This is called \emph{simultaneous Diophantine approximation}.

%%%%%%%%%%%%%%%%%%%%%%%%%%%%%%%%%%%%%%%%%%%%%%%%%%%%%%%%%%%%
\subsection{Simultaneous Dirichlet Theorem}
\label{subsec:simultaneous-dirichlet}
%%%%%%%%%%%%%%%%%%%%%%%%%%%%%%%%%%%%%%%%%%%%%%%%%%%%%%%%%%%%

A multidimensional pigeonhole argument gives:

\begin{theorem}[Simultaneous Dirichlet Approximation]
For real numbers

\[ \alpha_1,\ldots,\alpha_n \]

and a positive integer \(N\), there exists an integer

\[ 1\leq q\leq N^n \]

such that

\[ \boxed{\|q\alpha_i\|\leq\frac1N} \]

for every

\[ i=1,\ldots,n. \]
\end{theorem}

Thus several real numbers can be approximated simultaneously using the same
denominator.

%%%%%%%%%%%%%%%%%%%%%%%%%%%%%%%%%%%%%%%%%%%%%%%%%%%%%%%%%%%%
\subsection{Geometric Interpretation}
\label{subsec:diophantine-geometric-interpretation}
%%%%%%%%%%%%%%%%%%%%%%%%%%%%%%%%%%%%%%%%%%%%%%%%%%%%%%%%%%%%

The inequality

\[ |q\alpha-p|<\varepsilon \]

asks for integer lattice points

\[ (q,p)\in\Z^2 \]

lying close to the line

\[ y=\alpha x. \]

\begin{centereddiagram}
\begin{tikzpicture}[scale=0.72]
\draw[-{Latex[length=2.5mm]}] (0,0)--(8.5,0) node[right] {\(q\)};
\draw[-{Latex[length=2.5mm]}] (0,0)--(0,6.5) node[above] {\(p\)};

\foreach \x in {1,...,8}{
    \foreach \y in {1,...,6}{
        \fill (\x,\y) circle (1.2pt);
    }
}

\draw[thick,domain=0:8,samples=2] plot (\x,{0.68*\x})
node[right] {\(p=\alpha q\)};

\draw[dashed] (6,4)--(6,4.08);
\fill (6,4) circle (2.2pt);
\end{tikzpicture}
\end{centereddiagram}

Good rational approximations correspond to lattice points unusually close to
this line.

This viewpoint leads naturally to the geometry of numbers.

%%%%%%%%%%%%%%%%%%%%%%%%%%%%%%%%%%%%%%%%%%%%%%%%%%%%%%%%%%%%
\subsection{Geometry of Numbers Connection}
\label{subsec:geometry-of-numbers-connection}
%%%%%%%%%%%%%%%%%%%%%%%%%%%%%%%%%%%%%%%%%%%%%%%%%%%%%%%%%%%%

Instead of studying fractions directly, one can search for nonzero lattice
points inside carefully chosen geometric regions.

Minkowski's theorem and related lattice results provide powerful methods for
proving approximation theorems.

The conceptual translation is

\[
\boxed{
\text{rational approximation}
\longleftrightarrow
\text{lattice points near subspaces}.
}
\]

This geometric viewpoint generalizes naturally to higher dimensions.

%%%%%%%%%%%%%%%%%%%%%%%%%%%%%%%%%%%%%%%%%%%%%%%%%%%%%%%%%%%%
\subsection{Approximation Modulo One}
\label{subsec:approximation-modulo-one}
%%%%%%%%%%%%%%%%%%%%%%%%%%%%%%%%%%%%%%%%%%%%%%%%%%%%%%%%%%%%

The quantity

\[ \|q\alpha\| \]

also measures how close the fractional part of \(q\alpha\) is to \(0\).

Thus Diophantine approximation is closely connected to the sequence

\[ \{\alpha\},\{2\alpha\},\{3\alpha\},\ldots \]

inside the unit interval.

If \(\alpha\) is irrational, these fractional parts never repeat.

Their long-term distribution leads toward equidistribution theory.

%%%%%%%%%%%%%%%%%%%%%%%%%%%%%%%%%%%%%%%%%%%%%%%%%%%%%%%%%%%%
\subsection{Equidistribution Preview}
\label{subsec:equidistribution-preview}
%%%%%%%%%%%%%%%%%%%%%%%%%%%%%%%%%%%%%%%%%%%%%%%%%%%%%%%%%%%%

For irrational \(\alpha\), the sequence

\[ \{n\alpha\} \]

is equidistributed in

\[ [0,1). \]

Informally, this means that every interval receives its fair proportion of
the sequence.

For

\[ 0\leq a<b\leq1, \]

one has

\[
\frac{
\#\{1\leq n\leq N:\{n\alpha\}\in[a,b]\}
}{N}
\to b-a.
\]

This result connects approximation theory with harmonic analysis and
dynamical systems.

%%%%%%%%%%%%%%%%%%%%%%%%%%%%%%%%%%%%%%%%%%%%%%%%%%%%%%%%%%%%
\subsection{Approximation and Irrational Rotations}
\label{subsec:irrational-rotations}
%%%%%%%%%%%%%%%%%%%%%%%%%%%%%%%%%%%%%%%%%%%%%%%%%%%%%%%%%%%%

Consider moving around a circle by repeatedly adding an irrational angle
\(\alpha\) modulo \(1\).

The orbit is

\[ 0,\alpha,2\alpha,3\alpha,\ldots\pmod1. \]

\begin{centereddiagram}
\begin{tikzpicture}[scale=1.4]
\draw (0,0) circle (1.4);
\foreach \a in {0,47,94,141,188,235,282,329}{
    \fill ({1.4*cos(\a)},{1.4*sin(\a)}) circle (1.6pt);
}
\draw[-{Latex[length=2.5mm]}] (0,0)--({1.4*cos(47)},{1.4*sin(47)});
\node[right] at ({0.75*cos(47)},{0.75*sin(47)}) {\(\alpha\)};
\end{tikzpicture}
\end{centereddiagram}

Good rational approximations to \(\alpha\) correspond to times when the orbit
returns unusually close to its starting point.

%%%%%%%%%%%%%%%%%%%%%%%%%%%%%%%%%%%%%%%%%%%%%%%%%%%%%%%%%%%%
\subsection{Approximation and Numerical Computation}
\label{subsec:diophantine-numerical-computation}
%%%%%%%%%%%%%%%%%%%%%%%%%%%%%%%%%%%%%%%%%%%%%%%%%%%%%%%%%%%%

Rational approximation appears naturally in numerical computation.

A computer stores many quantities using finite representations.

Approximating a real number by

\[ \frac pq \]

can provide:

\begin{itemize}
    \item compact exact representations;
    \item efficient approximations to constants;
    \item stable rational reconstructions;
    \item detection of hidden algebraic structure;
    \item recovery of fractions from floating-point data.
\end{itemize}

Continued fractions are especially useful for reconstructing a rational
number from a decimal approximation.

%%%%%%%%%%%%%%%%%%%%%%%%%%%%%%%%%%%%%%%%%%%%%%%%%%%%%%%%%%%%
\subsection{Worked Example: Rational Reconstruction}
\label{subsec:worked-rational-reconstruction}
%%%%%%%%%%%%%%%%%%%%%%%%%%%%%%%%%%%%%%%%%%%%%%%%%%%%%%%%%%%%

\begin{workedexamplebox}{Recognizing a Decimal}

Suppose a computation produces

\[ x\approx0.4285714. \]

Identify a likely simple rational value.

\begin{tabularx}{\textwidth}{@{}>{\raggedright\arraybackslash}p{0.43\textwidth}X@{}}
Invert the decimal approximately. & \(1/x\approx2.333334\)\\
Separate the integer part. & \(2+\frac13\) approximately\\
Thus the continued fraction begins. & \([0;2,3]\)\\
Evaluate. & \([0;2,3]=\frac37\)
\end{tabularx}

\begin{answerbox}
\[ \boxed{x\approx\frac37.} \]
\end{answerbox}
\end{workedexamplebox}

%%%%%%%%%%%%%%%%%%%%%%%%%%%%%%%%%%%%%%%%%%%%%%%%%%%%%%%%%%%%
\subsection{Common Errors}
\label{subsec:diophantine-approximation-common-errors}
%%%%%%%%%%%%%%%%%%%%%%%%%%%%%%%%%%%%%%%%%%%%%%%%%%%%%%%%%%%%

\subsubsection{Ignoring the Denominator}

A very small error is not automatically impressive if the denominator is
enormous.

Approximation quality must be compared with \(q\).

%%%%%%%%%%%%%%%%%%%%%%%%%%%%%%%%%%%%%%%%%%%%%%%%%%%%%%%%%%%%
\subsubsection{Confusing Decimal Accuracy with Arithmetic Quality}

Two fractions may agree with a real number to the same number of decimal
places while using vastly different denominators.

The fraction with the smaller denominator may be arithmetically much more
efficient.

%%%%%%%%%%%%%%%%%%%%%%%%%%%%%%%%%%%%%%%%%%%%%%%%%%%%%%%%%%%%
\subsubsection{Assuming Every Good Approximation Is a Convergent}

Continued-fraction convergents provide canonical excellent approximations,
but other rational numbers can also approximate \(\alpha\) well.

%%%%%%%%%%%%%%%%%%%%%%%%%%%%%%%%%%%%%%%%%%%%%%%%%%%%%%%%%%%%
\subsubsection{Confusing Rational and Irrational Continued Fractions}

Rational numbers have finite simple continued fractions.

Irrational numbers have infinite simple continued fractions.

%%%%%%%%%%%%%%%%%%%%%%%%%%%%%%%%%%%%%%%%%%%%%%%%%%%%%%%%%%%%
\subsubsection{Confusing Periodic with Arbitrary Infinite Expansions}

Eventually periodic continued fractions characterize quadratic irrationals,
not all irrational numbers.

%%%%%%%%%%%%%%%%%%%%%%%%%%%%%%%%%%%%%%%%%%%%%%%%%%%%%%%%%%%%
\subsubsection{Assuming Large Irrationality Exponent Means Algebraic}

The opposite tendency occurs.

Algebraic irrational numbers have irrationality exponent \(2\), while
Liouville numbers have infinite irrationality exponent and are
transcendental.

%%%%%%%%%%%%%%%%%%%%%%%%%%%%%%%%%%%%%%%%%%%%%%%%%%%%%%%%%%%%
\subsubsection{Reversing Liouville's Implication}

Every Liouville number is transcendental.

Not every transcendental number is a Liouville number.

%%%%%%%%%%%%%%%%%%%%%%%%%%%%%%%%%%%%%%%%%%%%%%%%%%%%%%%%%%%%
\subsubsection{Misreading \(\|x\|\)}

In this section,

\[ \|x\| \]

usually means distance to the nearest integer, not absolute value and not a
vector norm.

%%%%%%%%%%%%%%%%%%%%%%%%%%%%%%%%%%%%%%%%%%%%%%%%%%%%%%%%%%%%
\subsubsection{Treating Dirichlet's Bound as the Best Bound for Every Number}

Dirichlet gives a universal guarantee.

Some numbers admit much stronger approximations.

Others, such as badly approximable numbers, resist substantial improvement.

%%%%%%%%%%%%%%%%%%%%%%%%%%%%%%%%%%%%%%%%%%%%%%%%%%%%%%%%%%%%
\subsubsection{Confusing Approximation with Equality}

A convergent such as

\[ \frac{99}{70} \]

is close to \(\sqrt2\), but it is not equal to \(\sqrt2\).

%%%%%%%%%%%%%%%%%%%%%%%%%%%%%%%%%%%%%%%%%%%%%%%%%%%%%%%%%%%%
\subsection{Applications and Connections}
\label{subsec:diophantine-approximation-applications}
%%%%%%%%%%%%%%%%%%%%%%%%%%%%%%%%%%%%%%%%%%%%%%%%%%%%%%%%%%%%

\begin{applicationbox}
\textbf{Elementary Number Theory.}
Continued fractions are closely connected to the Euclidean Algorithm,
greatest common divisors, and rational arithmetic.

\medskip

\textbf{Diophantine Equations.}
Approximations to algebraic numbers can produce integer solutions to equations
such as Pell equations.

\medskip

\textbf{Transcendental Number Theory.}
Liouville's theorem converts unusually strong rational approximation into a
method for proving transcendence.

\medskip

\textbf{Algebraic Number Theory.}
Approximation properties depend strongly on whether a number is rational,
algebraic irrational, or transcendental.

\medskip

\textbf{Geometry of Numbers.}
Rational approximation can be translated into the search for lattice points
near lines, planes, and higher-dimensional subspaces.

\medskip

\textbf{Dynamical Systems.}
The sequence \(n\alpha\pmod1\) describes an irrational rotation, connecting
approximation with recurrence and equidistribution.

\medskip

\textbf{Harmonic Analysis.}
Equidistribution can be studied using exponential sums and Fourier methods.

\medskip

\textbf{Numerical Mathematics.}
Continued fractions provide efficient rational approximations and rational
reconstruction algorithms.

\medskip

\textbf{Physics and Engineering.}
Rational approximation appears in resonance, frequency ratios, periodic
orbits, and near-commensurable systems.

\medskip

\textbf{Computer Science.}
Rational reconstruction and lattice-based approximation appear in symbolic
computation, cryptography, and algorithm design.
\end{applicationbox}

%%%%%%%%%%%%%%%%%%%%%%%%%%%%%%%%%%%%%%%%%%%%%%%%%%%%%%%%%%%%
\subsection{Exercises}
\label{subsec:diophantine-approximation-exercises}
%%%%%%%%%%%%%%%%%%%%%%%%%%%%%%%%%%%%%%%%%%%%%%%%%%%%%%%%%%%%

\begin{exercise}[1]
Compute

\[ \left|\sqrt2-\frac75\right|. \]
\end{exercise}

\begin{exercise}[1]
Compute

\[ \left|\pi-\frac{22}{7}\right|. \]
\end{exercise}

\begin{exercise}[1]
For

\[ \alpha=3.27,\qquad q=10, \]

choose the integer \(p\) nearest to \(q\alpha\) and determine the resulting
rational approximation.
\end{exercise}

\begin{exercise}[1]
Compute

\[ \|3.7\|,\qquad\|5.12\|,\qquad\|-2.8\|. \]
\end{exercise}

\begin{exercise}[2]
Show that

\[ \|x\|\leq\frac12 \]

for every real number \(x\).
\end{exercise}

\begin{exercise}[2]
Prove that

\[ \|x+n\|=\|x\| \]

for every \(n\in\Z\).
\end{exercise}

\begin{exercise}[2]
Explain why

\[ \left|\alpha-\frac pq\right|=\frac{|q\alpha-p|}{q}. \]
\end{exercise}

\begin{exercise}[2]
Use Dirichlet's approximation theorem with \(N=10\) to describe what can be
guaranteed for an arbitrary real number \(\alpha\).
\end{exercise}

\begin{exercise}[3]
Reproduce the pigeonhole proof of Dirichlet's approximation theorem.
\end{exercise}

\begin{exercise}[2]
Explain why Dirichlet's theorem implies infinitely many rational
approximations

\[ \left|\alpha-\frac pq\right|<\frac1{q^2} \]

when \(\alpha\) is irrational.
\end{exercise}

\begin{exercise}[1]
Convert

\[ \frac{23}{8} \]

to a finite simple continued fraction.
\end{exercise}

\begin{exercise}[1]
Convert

\[ \frac{47}{19} \]

to a finite simple continued fraction.
\end{exercise}

\begin{exercise}[2]
Evaluate

\[ [2;3,4]. \]
\end{exercise}

\begin{exercise}[2]
Evaluate

\[ [1;2,2,2]. \]
\end{exercise}

\begin{exercise}[2]
Find the first five convergents of

\[ [1;\overline2]. \]
\end{exercise}

\begin{exercise}[2]
Verify that

\[ \sqrt2=[1;\overline2]. \]
\end{exercise}

\begin{exercise}[2]
Compute the first six convergents of

\[ \varphi=[1;\overline1]. \]
\end{exercise}

\begin{exercise}[2]
Explain why the convergents of the golden ratio are ratios of consecutive
Fibonacci numbers.
\end{exercise}

\begin{exercise}[2]
Use the recurrence relations

\[ p_n=a_np_{n-1}+p_{n-2},\qquad q_n=a_nq_{n-1}+q_{n-2} \]

to compute the convergents of

\[ [3;7,15,1]. \]
\end{exercise}

\begin{exercise}[3]
Prove the identity

\[ p_nq_{n-1}-p_{n-1}q_n=(-1)^{n-1}. \]
\end{exercise}

\begin{exercise}[2]
Use the determinant identity to prove

\[ \gcd(p_n,q_n)=1. \]
\end{exercise}

\begin{exercise}[2]
Compare the errors of

\[ \frac{22}{7},\qquad\frac{333}{106},\qquad\frac{355}{113} \]

as approximations to \(\pi\).
\end{exercise}

\begin{exercise}[2]
Explain why a large continued-fraction partial quotient can produce an
exceptionally accurate convergent.
\end{exercise}

\begin{exercise}[2]
State Hurwitz's theorem and explain the significance of the constant

\[ \frac1{\sqrt5}. \]
\end{exercise}

\begin{exercise}[2]
Explain why the golden ratio is badly approximable.
\end{exercise}

\begin{exercise}[2]
Determine whether the continued fraction

\[ [2;1,3,1,3,1,3,\ldots] \]

is eventually periodic.
What does Lagrange's theorem imply about the represented number?
\end{exercise}

\begin{exercise}[2]
Verify that

\[ (x,y)=(3,2) \]

solves

\[ x^2-2y^2=1. \]
\end{exercise}

\begin{exercise}[2]
Verify that

\[ (x,y)=(17,12) \]

solves

\[ x^2-2y^2=1. \]
\end{exercise}

\begin{exercise}[3]
Explain conceptually why good rational approximations to \(\sqrt D\) are
related to Pell equations.
\end{exercise}

\begin{exercise}[2]
What lower bound does Dirichlet's theorem imply for the irrationality
exponent of every irrational real number?
\end{exercise}

\begin{exercise}[2]
State Liouville's approximation theorem in words.
\end{exercise}

\begin{exercise}[2]
Explain why every Liouville number is transcendental.
\end{exercise}

\begin{exercise}[2]
Explain why the converse

\[
\text{transcendental}
\Longrightarrow
\text{Liouville}
\]

is false.
\end{exercise}

\begin{exercise}[2]
Write the first four nonzero decimal positions of

\[ L=\sum_{n=1}^{\infty}10^{-n!}. \]
\end{exercise}

\begin{exercise}[3]
Show that truncations of Liouville's constant produce exceptionally strong
rational approximations.
\end{exercise}

\begin{exercise}[2]
State Roth's theorem conceptually.
\end{exercise}

\begin{exercise}[3]
Explain why Dirichlet's theorem and Roth's theorem together imply

\[ \mu(\alpha)=2 \]

for every algebraic irrational number \(\alpha\).
\end{exercise}

\begin{exercise}[2]
Draw the line

\[ y=\sqrt2\,x \]

and identify several integer lattice points lying relatively close to it.
Relate these points to rational approximations of \(\sqrt2\).
\end{exercise}

\begin{exercise}[2]
Explain geometrically why a lattice point \((q,p)\) close to

\[ y=\alpha x \]

produces a good rational approximation \(p/q\) to \(\alpha\).
\end{exercise}

\begin{exercise}[2]
Explain what it means for the sequence

\[ \{n\alpha\} \]

to be equidistributed modulo \(1\).
\end{exercise}

\begin{exercise}[3]
Explain the relationship between a good rational approximation

\[ \frac pq\approx\alpha \]

and a near return of the irrational rotation

\[ n\mapsto n\alpha\pmod1. \]
\end{exercise}

\begin{exercise}[2]
Use continued fractions to identify a simple rational approximation to

\[ 0.14285714. \]
\end{exercise}

\begin{exercise}[2]
Explain why continued fractions are useful for rational reconstruction in
numerical computation.
\end{exercise}

\begin{exercise}[3]
Compare the meanings of

\[ O\left(\frac1q\right),\qquad
\frac1{q^2},\qquad
\frac1{q^{2+\varepsilon}}. \]

Explain why the exponent matters in Diophantine approximation.
\end{exercise}

\begin{exercise}[3]
Explain the conceptual chain

\[
\text{Euclidean Algorithm}
\longrightarrow
\text{continued fraction}
\longrightarrow
\text{convergents}
\longrightarrow
\text{rational approximation}.
\]
\end{exercise}

\begin{exercise}[3]
Explain the conceptual chain

\[
\text{very strong rational approximation}
\longrightarrow
\text{Liouville's theorem}
\longrightarrow
\text{failure of algebraicity}
\longrightarrow
\text{transcendence}.
\]
\end{exercise}

%%%%%%%%%%%%%%%%%%%%%%%%%%%%%%%%%%%%%%%%%%%%%%%%%%%%%%%%%%%%
\subsection{Conceptual Summary}
\label{subsec:diophantine-approximation-summary}
%%%%%%%%%%%%%%%%%%%%%%%%%%%%%%%%%%%%%%%%%%%%%%%%%%%%%%%%%%%%

Diophantine approximation studies how efficiently real numbers can be
approximated by rational numbers.

An approximation

\[ \frac pq\approx\alpha \]

is measured using

\[ \left|\alpha-\frac pq\right|, \]

with particular attention to how the error depends on the denominator \(q\).

Dirichlet's approximation theorem guarantees that every irrational real
number has infinitely many rational approximations satisfying

\[ \boxed{\left|\alpha-\frac pq\right|<\frac1{q^2}.} \]

Continued fractions provide a systematic method for finding exceptionally
good rational approximations.

The convergents

\[ \frac{p_n}{q_n} \]

are generated recursively and satisfy strong best-approximation properties.

Rational numbers have finite continued fractions, while irrational numbers
have infinite continued fractions.

Lagrange's theorem gives the remarkable equivalence

\[
\boxed{
\text{quadratic irrational}
\iff
\text{eventually periodic continued fraction}.
}
\]

The golden ratio

\[ \varphi=[1;\overline1] \]

is a fundamental example of a badly approximable number and is closely
connected with Fibonacci numbers and the optimal constant in Hurwitz's
theorem.

The irrationality exponent

\[ \mu(\alpha) \]

measures how strongly \(\alpha\) can be approximated by rationals.

Dirichlet's theorem guarantees

\[ \mu(\alpha)\geq2 \]

for every irrational number.

Liouville's theorem shows that algebraic numbers cannot be approximated
arbitrarily well.

Liouville numbers violate every fixed algebraic approximation bound and are
therefore transcendental.

Roth's theorem gives the sharp result

\[ \boxed{\mu(\alpha)=2} \]

for every algebraic irrational number.

Simultaneous approximation extends these ideas to several real numbers.

Geometrically, rational approximation corresponds to finding lattice points
close to lines or higher-dimensional subspaces.

The central principle is

\[
\boxed{
\text{the arithmetic nature of a number is reflected in how well rationals
can approximate it}.
}
\]

%%%%%%%%%%%%%%%%%%%%%%%%%%%%%%%%%%%%%%%%%%%%%%%%%%%%%%%%%%%%
\subsection{Section Connections}
\label{subsec:diophantine-approximation-connections}
%%%%%%%%%%%%%%%%%%%%%%%%%%%%%%%%%%%%%%%%%%%%%%%%%%%%%%%%%%%%

\chapterconnection{
Diophantine Approximation connects rational arithmetic with the deeper
structure of irrational and transcendental numbers. The Euclidean Algorithm
leads naturally to continued fractions, whose convergents provide highly
efficient rational approximations. Quadratic irrational numbers are
characterized by periodic continued fractions and connect approximation
theory with Pell equations and algebraic number theory. Dirichlet's theorem
establishes the universal quadratic approximation scale, while Liouville's
and Roth's theorems show how algebraic structure restricts approximation.
These ideas lead directly into Transcendental Number Theory. Simultaneous
approximation and lattice interpretations connect the subject with Geometry
of Numbers, while fractional parts and irrational rotations connect it with
equidistribution, harmonic analysis, and dynamical systems. Computationally,
continued fractions provide practical algorithms for rational reconstruction
and numerical approximation.
}

%%%%%%%%%%%%%%%%%%%%%%%%%%%%%%%%%%%%%%%%%%%%%%%%%%%%%%%%%%%%
% SECTION SOURCE TRACKING
%%%%%%%%%%%%%%%%%%%%%%%%%%%%%%%%%%%%%%%%%%%%%%%%%%%%%%%%%%%%

%------------------------------------------------------------
% 2.6 Diophantine Approximation
%------------------------------------------------------------
%
% Source basis:
%   - Standard undergraduate and introductory graduate
%     Diophantine approximation theory.
%
% Used for:
%   - rational approximation and denominator-dependent error;
%   - nearest-integer distance notation;
%   - pigeonhole construction;
%   - Dirichlet's approximation theorem;
%   - continued fractions;
%   - convergents and recurrence relations;
%   - determinant identities;
%   - best-approximation properties;
%   - quadratic irrationals;
%   - periodic continued fractions;
%   - Lagrange's theorem;
%   - golden-ratio approximation;
%   - badly approximable numbers;
%   - Hurwitz's theorem;
%   - Pell-equation connections;
%   - irrationality exponents;
%   - Liouville's approximation theorem;
%   - Liouville numbers;
%   - Roth's theorem;
%   - simultaneous approximation;
%   - lattice interpretations;
%   - equidistribution preview;
%   - irrational rotations;
%   - rational reconstruction.
%
% Notes:
%   - Elementary continued-fraction arithmetic is developed
%     here because it is canonical to rational approximation.
%   - Full transcendence methods belong to
%     Transcendental Number Theory.
%   - Algebraic-number structure belongs to
%     Algebraic Number Theory.
%   - Lattice methods are developed more deeply in Geometry
%     and advanced Number Theory.
%   - Equidistribution and Fourier methods are developed
%     more deeply in Analysis.
%   - Probability and metric approximation are expanded in
%     later advanced material.
%
%%%%%%%%%%%%%%%%%%%%%%%%%%%%%%%%%%%%%%%%%%%%%%%%%%%%%%%%%%%%